\documentclass[frenchb, aspectratio=169]{beamer}

% Encodage, symboles et police
\usepackage[utf8]{inputenc}
\usepackage[T1]{fontenc}
\usepackage{lmodern}
\usepackage{eurosym}
\usepackage{relsize}
\usepackage{pifont}

\definecolor{itemcheck}{HTML}{357110}
\definecolor{itemcross}{HTML}{711110}

\newcommand*{\itemcheck}{\item[\textcolor{itemcheck}{\ding{51}}]}
\newcommand*{\itemcross}{\item[\textcolor{itemcross}{\ding{55}}]}

% Réglage des marges
\usepackage{geometry}

% Francisations
\usepackage{babel}

% Tables
\usepackage{booktabs}
\newcommand*{\tabtitle}[1]{\textbf{\textsc{\MakeLowercase{#1}}}}

% Graphismes
\usepackage{tikz}
\usetikzlibrary{scopes}
\usetikzlibrary{calc, fit}
\usetikzlibrary{positioning}
\usetikzlibrary{fadings}
\usetikzlibrary{shapes, shapes.misc}
\usetikzlibrary{decorations.pathreplacing}

\tikzset{
    font=\normalsize,
    % Permet de découvrir un graphisme TiKZ morceau par morceau
    invisible/.style={opacity=0},
    visible on/.style={alt={#1{}{invisible}}},
    alt/.code args={<#1>#2#3}{%
        \alt<#1>{\pgfkeysalso{#2}}{\pgfkeysalso{#3}}
    },
    % Style de flèche
    arrow/.style={thick, >=latex},
    % Style de flèche courte
    shortarrow/.style={
        arrow,
        shorten >=10pt,
        shorten <=10pt
    },
    % Style d’image
    image/.style={align=center},
}

% Liens
\usepackage{hyperref}

% Figures
\usepackage{graphicx}
\usepackage[justification=justified]{caption}
\usepackage{subcaption}
\graphicspath{{./images/}}

% Titre, date et auteurs
\newcommand{\et}{\hspace{.3em}{\smaller et}\hspace{.5em}}

\title{Scrum}
\subtitle{Une méthode Agile}
\author{Rémi Cérès \et Mattéo Delabre}
\institute{HLIN408 -- Techniques de communication\\Université de Montpellier}
\titlegraphic{
    \includegraphics[scale=.25]{images/logo-um}\quad
    \includegraphics[scale=.125]{images/logo-fds}
}
\date{24 avril 2017}

% Thème
\usetheme{Scrum}

\begin{document}

% Page de titre
\begin{frame}[plain]
\titlepage
\end{frame}

% Page de définition de Scrum
\begin{frame}{}
    \vfill
    \centering
    \begin{tikzpicture}
        \node[
            thick, draw=fg,
            inner sep=1cm,
            align=center,
        ]
        (frame)
        {\alert<2->{Méthode Agile} de\\\alert<2->{gestion de projet} informatique.};

        \node[
            fill=bg,
            minimum width=5cm
        ] at (frame.north) {\usebeamerfont{title}scrum};
    \end{tikzpicture}
    \vfill
\end{frame}

% Table des matières
\begin{frame}{Vue d’ensemble}
    \tableofcontents
\end{frame}

% Contenu
\section{Qu’est-ce que la gestion de projet ?}

\begin{frame}[t]{Qu’est-ce qu’un projet ?}
    \textbf{Projet :} processus permettant d’atteindre un objectif fixé.

    \vspace{1.4cm}

    \begin{center}
    \begin{tikzpicture}
        \node[image] (idea) {\includegraphics[width=1cm]{icon-idea}\\Idée};
        \node[image] (project) at (8cm, 0) {\includegraphics[width=1cm]{icon-target}\\Objectif};

        \draw[->, shortarrow] (idea) -- node[below] {\textit{Projet}} (project);
    \end{tikzpicture}
    \end{center}

    \note{note test}
\end{frame}

\begin{frame}[t]{Comment gérer un projet ?}
    \textbf{Gestion de projet :} outils et méthodes permettant d’organiser un projet.

    \vfill

    \begin{tikzpicture}
        \node[image] (project) {\includegraphics[width=1cm]{icon-project}\\Projet};
        \node[image] (team) at (\textwidth - 2cm, 0) {\includegraphics[width=1cm]{icon-team}\\Équipe};

        {[visible on=<2->]
            \node[image, right=2cm of project] (task2) {\includegraphics[width=1cm]{icon-task}};
            \node[image, above=.5cm of task2] (task1) {\includegraphics[width=1cm]{icon-task}};
            \node[image, below=.5cm of task2] (task3) {\includegraphics[width=1cm]{icon-task}\\Tâches};

            \draw[->, shortarrow] (project) edge[out=20, in=200] (task1);
            \draw[->, shortarrow] (project) -- (task2);
            \draw[->, shortarrow] (project) edge[out=-20, in=160] (task3);
        }

        {[visible on=<3->]
            \node[image, left=2cm of team] (role2) {\includegraphics[width=1cm]{icon-man}};
            \node[image, above=.5cm of role2] (role1) {\includegraphics[width=1cm]{icon-woman}};
            \node[image, below=.5cm of role2] (role3) {\includegraphics[width=1cm]{icon-woman}\\Rôles};

            \draw[->, shortarrow] (team) edge[out=160, in=-20] (role1);
            \draw[->, shortarrow] (team) -- (role2);
            \draw[->, shortarrow] (team) edge[out=200, in=20] (role3);
        }

        {[visible on=<4->]
            \draw[<->, shortarrow] (task1) -- (role1);
            \draw[<->, shortarrow] (task1) -- (role3);
            \draw[<->, shortarrow] (task2) -- (role1);
            \draw[<->, shortarrow] (task3) -- (role2);
        }
    \end{tikzpicture}
\end{frame}

\newcommand{\imgcaption}[2]{%
    \begin{minipage}[t]{0.135\linewidth}%
        \centering
        \includegraphics[width=\linewidth]{#1}\\%
        {\scriptsize #2}%
    \end{minipage}%
}

\begin{frame}{Pourquoi utiliser une méthode ?}
\makebox[\linewidth][c]{%
    \begin{minipage}{\dimexpr\textwidth+4em\relax}
        \imgcaption{images/arbre-1}{Comment le client a exprimé ses besoins}%
        \hfill%
        \imgcaption{images/arbre-2}{Comment le chef de projet les a compris}%
        \hfill%
        \imgcaption{images/arbre-3}{Ce que les ingénieurs ont conçu}%
        \hfill%
        \imgcaption{images/arbre-4}{Ce que les programmeurs ont codé}%
        \hfill%
        \imgcaption{images/arbre-5}{Comment le projet a été vendu}%
        \hfill%
        \imgcaption{images/arbre-6}{Ce qu’à reçu le client}%
        \hfill%
        \imgcaption{images/arbre-7}{Ce que voulait vraiment le client}%
    \end{minipage}%
}
\end{frame}

\section{Une méthode classique : le cycle en V}

\begin{frame}{Point d’histoire}
    \begin{itemize}
        \item Mis au point par les départements de défense étasunien et allemand.
        \item Initialement destiné aux systèmes de satellites.
        \item Adapté à l’ingénierie logicielle dans les années 1980.
    \end{itemize}
\end{frame}

\begin{frame}{Comment fonctionne le cycle en V ?}
    \tikzset{
        phasebox/.style={
            rectangle,
            thick,
            align=center,
            draw=fg,
            minimum height=.75cm
        }
    }

    \begin{tikzpicture}
        {[
            node distance=.5cm,
            every node/.style={phasebox},
            edge from parent/.style={},
            level distance=1.6cm,
            level 1/.style={sibling distance=2.5cm},
            level 2/.style={sibling distance=1.5cm}
        ]
            % Arbre en V
            \node [visible on=<4->] (codage) {Programmation} [grow'=up]
            child {
                node [visible on=<3->] (conc-det) {Conception\\modulaire}
                child {
                    node [visible on=<2->](conc-arc) {Conception\\architecturale}
                    child {
                        node (anal-bes) {Analyse des\\besoins}
                        child[missing]
                    }
                    child[missing]
                }
                child[missing]
            }
            child {
                node [visible on=<5->] (test-unit) {Test\\unitaire}
                child[missing]
                child {
                    node [visible on=<6->] (test-int) {Test\\d’intégration}
                    child[missing]
                    child {
                        node [visible on=<7->] (test-acc) {Test\\d’acceptation}
                    }
                }
            };
        }

        % Flèches ordonnées
        \draw[->, arrow, visible on=<2->] (anal-bes) -- (conc-arc);
        \draw[->, arrow, visible on=<3->] (conc-arc) -- (conc-det);
        \draw[->, arrow, visible on=<4->] (conc-det) -- (codage);
        \draw[->, arrow, visible on=<5->] (codage) -- (test-unit);
        \draw[->, arrow, visible on=<6->] (test-unit) -- (test-int);
        \draw[->, arrow, visible on=<7->] (test-int) -- (test-acc);

        % Flèches de retour
        \draw[->, arrow, visible on=<5->] (test-unit) -- (conc-det);
        \draw[->, arrow, visible on=<6->] (test-int) -- (conc-arc);
        \draw[->, arrow, visible on=<7->] (test-acc) -- (anal-bes);

        % Rôles
        \node[image, right=of test-acc] (role-user)
            {\includegraphics[width=1cm]{icon-user}\\Utilisateur};

        \node[image, visible on=<4->] (role-coding) at
            (codage -| role-user)
            {\includegraphics[width=1cm]{icon-coding}\\Programmeurs};

        \node[image, visible on=<3->] (role-engineer) at
            ($(role-coding)!0.5!(role-user)$)
            {\includegraphics[width=1cm]{icon-engineer}\\Ingénieurs};

        \draw[
            decoration={brace, raise=1.75cm},
            thick, decorate, visible on=<8->
        ] (role-user.north) --
            node[image, right=2.25cm]
            {\includegraphics[width=1cm]{icon-boss}\\Chef de\\projet}
          (role-coding.south);
    \end{tikzpicture}
\end{frame}

\begin{frame}{Avantages et inconvénients}
    \begin{itemize}
        \itemcheck Simple à comprendre et à appréhender.
        \itemcheck Standard historique de l’industrie.
        \itemcheck Fixe un accord avec le client sur les spécifications.
        \itemcheck Intègre les tests au processus de création.
    \end{itemize}
    \pause
    \begin{itemize}
        \itemcross Crée une hiérarchie importante.
        \itemcross Déconnexion entre le client et la création du produit.
        \itemcross Désynchronisation entre conception et implantation.
        \itemcross Pas suffisamment flexible (changement de besoin du client).
    \end{itemize}
\end{frame}

\section{Une méthode Agile ?}

\begin{frame}{Genèse}{Développement d’alternatives « légères »}
    \begin{tabular}{llr}
        \tabtitle{Nom de la méthode} & \tabtitle{Auteur} & \tabtitle{Année}\\
        \toprule
        Rapid Application Development & James Martin & 1991\\
        \midrule
        Dynamic Systems Development Method & --- & 1994\\
        \midrule
        \textbf{Scrum} & \textbf{Ken Schwaber} & \textbf{1995}\\
        \midrule
        Crystal Clear & Alistair Cockburn & 1996\\
        \midrule
        Extreme Programming & Kent Beck & 1996\\
        \bottomrule
    \end{tabular}
\end{frame}

\newcommand*{\itemdate}{\raisebox{-0.2\height}{\includegraphics[height=.5cm]{icon-date}}}
\newcommand*{\itemplace}{\raisebox{-0.2\height}{\includegraphics[height=.5cm]{icon-place}}}
\newcommand*{\itemgoal}{\raisebox{-0.2\height}{\includegraphics[height=.5cm]{icon-target}}}

\begin{frame}{Genèse}{Réunion pour les méthodes « légères »}
    17 représentants des méthodes « légères » se réunissent.

    \vspace{1em}
    \begin{itemize}
        \item[\itemdate] \textbf{Date :} février 2001, pendant 2 jours.
        \item[\itemplace] \textbf{Lieu :} station de ski de l'Utah aux États-Unis.
        \item[\itemgoal] \textbf{Objectif :} discuter de leur expérience, mettre en commun leurs idées.
    \end{itemize}
\end{frame}

{
    \usebackgroundtemplate{%
        \includegraphics[width=\paperwidth]{manifeste}%
    }

    \begin{frame}{Manifeste Agile}{Définition du terme « Agile »}
        Quatre valeurs fondamentales forment le socle commun des méthodes Agiles.

        \vspace{1em}
        \begin{tabular}{rcl}
            \textbf{Humains et interactions}&$\geqslant$&processus et outils\\
            \textbf{Logiciels opérationnels}&$\geqslant$&documentation exhaustive\\
            \textbf{Collaboration avec les clients}&$\geqslant$&négociation contractuelle\\
            \textbf{Adaptation au changement}&$\geqslant$&suivi d’un plan
        \end{tabular}
    \end{frame}
}

\begin{frame}{Manifeste Agile}{Conséquences et rayonnement philosophique}
    \begin{itemize}
        \item Marque la formation de l'Agile Alliance.
        \item Renforce l'unité des méthodes agiles.
        \item Place les méthodes Agiles comme alternatives crédibles.
    \end{itemize}
\end{frame}

\section{Le cadre de travail Scrum}

\begin{frame}{Philosophie générale de Scrum}
    \begin{columns}
        \begin{column}{.55\textwidth}
            \begin{itemize}
                \item Scrum (« mêlée » en français) s'inspire des valeurs collectives du rugby.
                \item Approche incrémentale et itérative.
                \item Travail en équipe au cœur du projet.
            \end{itemize}
        \end{column}
        \begin{column}{.45\textwidth}
            \includegraphics[width=\linewidth]{rugby}
            \vfill
            \hfill{\scriptsize \textbf{Source :} Kiwi Flickr (CC BY 2.0)}\\[-.3em]
            \hfill{\scriptsize\url{https://flic.kr/p/3bnXZ}}
        \end{column}
    \end{columns}
\end{frame}

\definecolor{scrumrelease}{HTML}{005480}
\definecolor{scrumsprint}{HTML}{800014}
\definecolor{scrummilestone}{HTML}{806C00}

\tikzset{
    baseblock/.style={
        minimum height=3ex,
        rounded corners,
        fill=fg,
        text=white,
        inner sep=1.2ex,
    },
    block/.style={
        baseblock,
        text depth=0ex,
        text height=1.2ex,
    },
    release/.style={
        block,
        fill=scrumrelease,
    },
    sprint/.style={
        block,
        fill=scrumsprint,
    },
    point/.style={
        fill=fg,
        circle,
        inner sep=3pt,
    }
}

\begin{frame}[t]{Cadre temporel de Scrum}{Sprint : unité de base}
    \vspace{1.5cm}
    \centering\tikz\node[sprint] {Sprint};

    \begin{itemize}
        \item \textbf{Objectif : apporter un incrément.}
        \item Durée fixée à l’avance en semaines (2 -- 3 semaines).
        \item Échéance non modifiable : pas de retard, budget fixé.
        \item L’équipe reste la même.
    \end{itemize}
\end{frame}

\begin{frame}[t]{Cadre temporel de Scrum}{Déroulement d'un sprint}
    \vspace{1.5cm}
    \begin{tikzpicture}
        \fill[color=scrumsprint!50]
            (.5cm, 0) {[rounded corners] --
            (.5cm, 2.15) --
            (\linewidth - 2pt - 1cm, 2.15)} --
            (\linewidth - 2pt - 1cm, 0);
        \draw[arrow, very thick, ->] (0, 0) -- (\linewidth - 2pt, 0);

        \node[sprint] (sprint-master) at (.5\linewidth - .5pt, 2.15) {Sprint};

        \node[block, anchor=west] (sprint-plan) at (.5, 0) {Planification};
        \node[block, anchor=east, align=center, visible on=<4->] (sprint-retro) at (\linewidth - 2pt - 1cm, 0) {Rétrospective};
        \node[baseblock, left=.6cm of sprint-retro, align=center, visible on=<3->] (sprint-revue) {Revue de\\sprint};
        \node[baseblock, align=center, right=.6cm of sprint-plan, visible on=<2->] (sprint-daily) {Mêlées\\quotidiennes};

        \draw[arrow, very thick, ->, visible on=<2->] (sprint-daily) edge[in=-110,out=-70,loop,looseness=4] (sprint-daily);
    \end{tikzpicture}
\end{frame}

\begin{frame}[t]{Cadre temporel de Scrum}{Release : agrégat de sprints}
    \vspace{1.5cm}
    \centering\tikz\node[release] {Release};

    \begin{itemize}
        \item \textbf{Objectif : livrer une version.}
        \item Durée flexible (3 -- 4 mois).
    \end{itemize}
\end{frame}

\begin{frame}[t]{Cadre temporel de Scrum}{Déroulement d'une release}
    \vspace{1.5cm}
    \begin{tikzpicture}
        \fill[color=scrumrelease!50]
            (.5cm, 0) {[rounded corners] --
            (.5cm, 2) --
            (\linewidth - 2pt - 1cm, 2)} --
            (\linewidth - 2pt - 1cm, 0);
        \draw[arrow, very thick, ->] (0, 0) -- (\linewidth - 2pt, 0);

        \node[release] (release-master) at (.5\linewidth - .5pt, 2) {Release};

        \node[block, anchor=west] (release-prep) at (.5, 0) {Préparation};
        \node[sprint, visible on=<2->] (release-sprint) at (.5\linewidth, 0) {Sprints};
        \node[block, anchor=east, visible on=<3->] (release-final) at (\linewidth - 2pt - 1cm, 0) {Finalisation};

        \draw[arrow, very thick, ->, visible on=<2->] (release-sprint) edge[in=-120,out=-60,loop,looseness=5] (release-sprint);
    \end{tikzpicture}
\end{frame}

\begin{frame}{Cadre temporel de Scrum}{Déroulement d'un projet}
    \pgfdeclarelayer{background}
    \pgfdeclarelayer{front}
    \pgfsetlayers{background,front}

    \begin{tikzpicture}
        \begin{pgfonlayer}{front}
            \node[point] (start) {};
            \node[below=0cm of start, align=center] {Début du\\projet};

            \node[release, right=1cm of start] (release1) {\color{white}{Release 1}};
            \node[release, right=1cm of release1] (release2) {\color{white}{Release 2}};
            \node[release, right=1cm of release2] (release3) {\color{white}{Release 3}};

            \draw[
                decoration={brace, raise=.5cm},
                thick, decorate
            ] (release2.east) --
                node[below=.75cm] {3 – 4 mois}
                (release2.west);

            \node[point] (end) at (\linewidth - 2cm, 0) {};
            \node[below=0cm of end, align=center] {Fin du\\projet};
        \end{pgfonlayer}

        \begin{pgfonlayer}{background}
            \draw[very thick] (start.center) -- (release3.west);
            \draw[arrow, very thick, ->, dash pattern=on 6pt off 6pt]
                (release3.west) -- (end);
        \end{pgfonlayer}
    \end{tikzpicture}
    \vfill
    \begin{itemize}
        \item À la différence des sprints, les releases sont espacées.
        \item Permet de faire une rétrospective et de faire la fête.
    \end{itemize}
\end{frame}

\begin{frame}{Cadre temporel de Scrum}{En bref}
    \begin{itemize}
        \item Projet découpé en \tikz[baseline=(rel.base)] \node[release, scale=.7] (rel) {Releases}; $\rightarrow$ versions livrables.
        \item Releases découpées en \tikz[baseline=(spr.base)] \node[sprint, scale=.7] (spr) {Sprints}; $\rightarrow$ incréments au produit.
        \item Réunions régulières pour affiner le cycle de développement du produit.
        \item Coordination avec le client après chaque sprint.
    \end{itemize}
\end{frame}

\begin{frame}{Cadre humain de Scrum}{L'équipe selon Scrum}
    \textbf{Fonctionnement}

    \begin{itemize}
        \item L'équipe a le pouvoir sur la réalisation du projet.
        \item Équipe non-hiérarchisée, auto-organisée et indépendante.
    \end{itemize}

    \pause
    \vfill
    \textbf{Critères de choix de l'équipe}

    \begin{itemize}
        \item Taille de l'équipe : 2 à 9 personnes (optimum à 7).
        \item Équipe pluridisciplinaire.
        \item Cohésion et proximité géographique.
    \end{itemize}
\end{frame}

\begin{frame}{Cadre humain de Scrum}{Vue d'ensemble des rôles}
    % Schéma des différents rôles occupés par les acteurs de Scrum

\begin{tikzpicture}
    % L'équipe : ScrumMaster et programmeurs
    \foreach \prognum in {1,2,3,4,5}{
        \node[image, inner sep=0, outer sep=0] (prog-\prognum) at (\prognum*360/5+18 : 2.4cm and 1.8cm)
            {\includegraphics[width=1cm]{icon-coding}};
    }

    \draw ($(prog-3.south)!0.5!(prog-4.south)$)
        node[below] (prog-text) {Programmeurs};

    \node[image] (sm) {\includegraphics[width=1cm]{icon-trainer}\\ScrumMaster};

    \node[
        draw=fg, thick, ellipse,
        inner xsep=0cm, inner ysep=-.6cm,
        fit=(prog-1) (prog-2) (prog-3) (prog-4) (prog-5) (prog-text) (sm)
    ] (team) {}
        node[above] at (team.north) {\scshape\bfserieséquipe};

    % Interface du PO
    \node[
        rectangle, fill=white,
        yshift=-.25cm, inner xsep=1cm, inner ysep=1.5cm
    ] at (team.east) {};

    \node[image] (po) at (sm-|team.east)
        {\includegraphics[width=1cm]{icon-user}}
        node[anchor=north, align=center] at (po.south) {Product\\Owner};

    % Parties prenantes
    \node[image, right=of po] (pp2)
        {\includegraphics[width=1cm]{icon-woman}};

    \node[image, above=0cm of pp2] (pp1)
        {\includegraphics[width=1cm]{icon-man}};

    \node[image, below=0cm of pp2] (pp3)
        {\includegraphics[width=1cm]{icon-man}\\Parties\\prenantes};

    % Experts
    \node [image, left=of sm-|team.west] (exp2)
        {\includegraphics[width=1cm]{icon-scientist}};

    \node [image, above=0cm of exp2] (exp1)
        {\includegraphics[width=1cm]{icon-scientist}};

    \node [image, below=0cm of exp2] (exp3)
        {\includegraphics[width=1cm]{icon-scientist}\\Experts};
\end{tikzpicture}

\end{frame}

\begin{frame}{Cadre humain de Scrum}
\begin{columns}
    \begin{column}{.4\textwidth}
        \hspace{2em}\textbf{Scrum Master}
        \begin{itemize}
            \item Coach de l'équipe.
            \item Au service de l'équipe\\(pas un supérieur).
            \item Garant de l'application\\de Scrum.
            \item Aide à éliminer les obstacles.
        \end{itemize}
    \end{column}%
    \begin{column}{.6\textwidth}
        \newcommand{\scrumfocus}{sm}
        % Schéma des différents rôles occupés par les acteurs de Scrum,
% sans les labels pour prendre moins de place

\newcommand{\opacityat}[2]{
    \ifthenelse{\equal{\scrumfocus}{#1}}%
        {\begin{scope}[opacity=1, transparency group]#2\end{scope}}%
        {\begin{scope}[opacity=.3, transparency group]#2\end{scope}}
}

\begin{tikzpicture}
    % L'équipe : ScrumMaster et programmeurs
    \foreach \prognum in {1,2,3,4,5}{
        \opacityat{prog}{
            \node[image, inner sep=0, outer sep=0]
                (prog-\prognum) at (\prognum*360/5+18 : 1.4cm and 1.2cm)
                {\includegraphics[width=.8cm]{icon-coding}};
        }
    }

    \opacityat{sm}{
        \node[image] (sm) {\includegraphics[width=.8cm]{icon-trainer}};
    }

    \node[
        draw=fg, thick, ellipse,
        inner sep=0cm,
        fit=(prog-1) (prog-2) (prog-3) (prog-4) (prog-5) (sm)
    ] (team) {};

    % Interface du PO
    \node[
        rectangle, fill=white,
        inner xsep=.6cm, inner ysep=.6cm
    ] at (sm-|team.east) {};

    \opacityat{po}{
        \node[image] (po) at (sm-|team.east)
            {\includegraphics[width=.8cm]{icon-user}};
    }

    % Parties prenantes
    \opacityat{pp}{
        \node[image, right=.2cm of po] (pp2)
            {\includegraphics[width=.8cm]{icon-woman}};

        \node[image, above=0cm of pp2] (pp1)
            {\includegraphics[width=.8cm]{icon-man}};

        \node[image, below=0cm of pp2] (pp3)
            {\includegraphics[width=.8cm]{icon-man}};
    }

    % Experts
    \opacityat{exp}{
        \node [image, left=.2cm of sm-|team.west] (exp2)
            {\includegraphics[width=.8cm]{icon-scientist}};

        \node [image, above=0cm of exp2] (exp1)
            {\includegraphics[width=.8cm]{icon-scientist}};

        \node [image, below=0cm of exp2] (exp3)
            {\includegraphics[width=.8cm]{icon-scientist}};
    }
\end{tikzpicture}

    \end{column}
\end{columns}
\end{frame}

\begin{frame}{Cadre humain de Scrum}
\begin{columns}
    \begin{column}{.4\textwidth}
        \hspace{2em}\textbf{Product Owner}
        \begin{itemize}
            \item Médiateur entre l'équipe et les parties prenantes.
            \item Vision globale du produit (fonctionnalités, attentes, objectifs).
            \item Gère et priorise les tâches et répartit le temps.
        \end{itemize}
    \end{column}%
    \begin{column}{.6\textwidth}
        \newcommand{\scrumfocus}{po}
        % Schéma des différents rôles occupés par les acteurs de Scrum,
% sans les labels pour prendre moins de place

\newcommand{\opacityat}[2]{
    \ifthenelse{\equal{\scrumfocus}{#1}}%
        {\begin{scope}[opacity=1, transparency group]#2\end{scope}}%
        {\begin{scope}[opacity=.3, transparency group]#2\end{scope}}
}

\begin{tikzpicture}
    % L'équipe : ScrumMaster et programmeurs
    \foreach \prognum in {1,2,3,4,5}{
        \opacityat{prog}{
            \node[image, inner sep=0, outer sep=0]
                (prog-\prognum) at (\prognum*360/5+18 : 1.4cm and 1.2cm)
                {\includegraphics[width=.8cm]{icon-coding}};
        }
    }

    \opacityat{sm}{
        \node[image] (sm) {\includegraphics[width=.8cm]{icon-trainer}};
    }

    \node[
        draw=fg, thick, ellipse,
        inner sep=0cm,
        fit=(prog-1) (prog-2) (prog-3) (prog-4) (prog-5) (sm)
    ] (team) {};

    % Interface du PO
    \node[
        rectangle, fill=white,
        inner xsep=.6cm, inner ysep=.6cm
    ] at (sm-|team.east) {};

    \opacityat{po}{
        \node[image] (po) at (sm-|team.east)
            {\includegraphics[width=.8cm]{icon-user}};
    }

    % Parties prenantes
    \opacityat{pp}{
        \node[image, right=.2cm of po] (pp2)
            {\includegraphics[width=.8cm]{icon-woman}};

        \node[image, above=0cm of pp2] (pp1)
            {\includegraphics[width=.8cm]{icon-man}};

        \node[image, below=0cm of pp2] (pp3)
            {\includegraphics[width=.8cm]{icon-man}};
    }

    % Experts
    \opacityat{exp}{
        \node [image, left=.2cm of sm-|team.west] (exp2)
            {\includegraphics[width=.8cm]{icon-scientist}};

        \node [image, above=0cm of exp2] (exp1)
            {\includegraphics[width=.8cm]{icon-scientist}};

        \node [image, below=0cm of exp2] (exp3)
            {\includegraphics[width=.8cm]{icon-scientist}};
    }
\end{tikzpicture}

    \end{column}
\end{columns}
\end{frame}

\begin{frame}{Cadre humain de Scrum}
\begin{columns}
    \begin{column}{.4\textwidth}
        \hspace{2em}\textbf{Programmeurs}
        \begin{itemize}
            \item Effectuent les tâches du projet.
            \item Partagent leurs compétences et les mettent au service de l'équipe.
            \item Participent à la prise de décisions.
        \end{itemize}
    \end{column}%
    \begin{column}{.6\textwidth}
        \newcommand{\scrumfocus}{prog}
        % Schéma des différents rôles occupés par les acteurs de Scrum,
% sans les labels pour prendre moins de place

\newcommand{\opacityat}[2]{
    \ifthenelse{\equal{\scrumfocus}{#1}}%
        {\begin{scope}[opacity=1, transparency group]#2\end{scope}}%
        {\begin{scope}[opacity=.3, transparency group]#2\end{scope}}
}

\begin{tikzpicture}
    % L'équipe : ScrumMaster et programmeurs
    \foreach \prognum in {1,2,3,4,5}{
        \opacityat{prog}{
            \node[image, inner sep=0, outer sep=0]
                (prog-\prognum) at (\prognum*360/5+18 : 1.4cm and 1.2cm)
                {\includegraphics[width=.8cm]{icon-coding}};
        }
    }

    \opacityat{sm}{
        \node[image] (sm) {\includegraphics[width=.8cm]{icon-trainer}};
    }

    \node[
        draw=fg, thick, ellipse,
        inner sep=0cm,
        fit=(prog-1) (prog-2) (prog-3) (prog-4) (prog-5) (sm)
    ] (team) {};

    % Interface du PO
    \node[
        rectangle, fill=white,
        inner xsep=.6cm, inner ysep=.6cm
    ] at (sm-|team.east) {};

    \opacityat{po}{
        \node[image] (po) at (sm-|team.east)
            {\includegraphics[width=.8cm]{icon-user}};
    }

    % Parties prenantes
    \opacityat{pp}{
        \node[image, right=.2cm of po] (pp2)
            {\includegraphics[width=.8cm]{icon-woman}};

        \node[image, above=0cm of pp2] (pp1)
            {\includegraphics[width=.8cm]{icon-man}};

        \node[image, below=0cm of pp2] (pp3)
            {\includegraphics[width=.8cm]{icon-man}};
    }

    % Experts
    \opacityat{exp}{
        \node [image, left=.2cm of sm-|team.west] (exp2)
            {\includegraphics[width=.8cm]{icon-scientist}};

        \node [image, above=0cm of exp2] (exp1)
            {\includegraphics[width=.8cm]{icon-scientist}};

        \node [image, below=0cm of exp2] (exp3)
            {\includegraphics[width=.8cm]{icon-scientist}};
    }
\end{tikzpicture}

    \end{column}
\end{columns}
\end{frame}

\begin{frame}{Cadre humain de Scrum}
\begin{columns}
    \begin{column}{.4\textwidth}
        \hspace{2em}\textbf{Parties prenantes}
        \begin{itemize}
            \item Personnes intéressées par les résultats (clients, utilisateurs, managers).
            \item Participent aux revues.
        \end{itemize}
    \end{column}%
    \begin{column}{.6\textwidth}
        \newcommand{\scrumfocus}{pp}
        % Schéma des différents rôles occupés par les acteurs de Scrum,
% sans les labels pour prendre moins de place

\newcommand{\opacityat}[2]{
    \ifthenelse{\equal{\scrumfocus}{#1}}%
        {\begin{scope}[opacity=1, transparency group]#2\end{scope}}%
        {\begin{scope}[opacity=.3, transparency group]#2\end{scope}}
}

\begin{tikzpicture}
    % L'équipe : ScrumMaster et programmeurs
    \foreach \prognum in {1,2,3,4,5}{
        \opacityat{prog}{
            \node[image, inner sep=0, outer sep=0]
                (prog-\prognum) at (\prognum*360/5+18 : 1.4cm and 1.2cm)
                {\includegraphics[width=.8cm]{icon-coding}};
        }
    }

    \opacityat{sm}{
        \node[image] (sm) {\includegraphics[width=.8cm]{icon-trainer}};
    }

    \node[
        draw=fg, thick, ellipse,
        inner sep=0cm,
        fit=(prog-1) (prog-2) (prog-3) (prog-4) (prog-5) (sm)
    ] (team) {};

    % Interface du PO
    \node[
        rectangle, fill=white,
        inner xsep=.6cm, inner ysep=.6cm
    ] at (sm-|team.east) {};

    \opacityat{po}{
        \node[image] (po) at (sm-|team.east)
            {\includegraphics[width=.8cm]{icon-user}};
    }

    % Parties prenantes
    \opacityat{pp}{
        \node[image, right=.2cm of po] (pp2)
            {\includegraphics[width=.8cm]{icon-woman}};

        \node[image, above=0cm of pp2] (pp1)
            {\includegraphics[width=.8cm]{icon-man}};

        \node[image, below=0cm of pp2] (pp3)
            {\includegraphics[width=.8cm]{icon-man}};
    }

    % Experts
    \opacityat{exp}{
        \node [image, left=.2cm of sm-|team.west] (exp2)
            {\includegraphics[width=.8cm]{icon-scientist}};

        \node [image, above=0cm of exp2] (exp1)
            {\includegraphics[width=.8cm]{icon-scientist}};

        \node [image, below=0cm of exp2] (exp3)
            {\includegraphics[width=.8cm]{icon-scientist}};
    }
\end{tikzpicture}

    \end{column}
\end{columns}
\end{frame}

\begin{frame}{Cadre humain de Scrum}
\begin{columns}
    \begin{column}{.4\textwidth}
        \hspace{2em}\textbf{Experts}
        \begin{itemize}
            \item Jokers utilisés ponctuellement.
            \item Compétents dans un domaine particulier.
            \item Permettent d'outrepasser des difficultés.
        \end{itemize}
    \end{column}%
    \begin{column}{.6\textwidth}
        \newcommand{\scrumfocus}{exp}
        % Schéma des différents rôles occupés par les acteurs de Scrum,
% sans les labels pour prendre moins de place

\newcommand{\opacityat}[2]{
    \ifthenelse{\equal{\scrumfocus}{#1}}%
        {\begin{scope}[opacity=1, transparency group]#2\end{scope}}%
        {\begin{scope}[opacity=.3, transparency group]#2\end{scope}}
}

\begin{tikzpicture}
    % L'équipe : ScrumMaster et programmeurs
    \foreach \prognum in {1,2,3,4,5}{
        \opacityat{prog}{
            \node[image, inner sep=0, outer sep=0]
                (prog-\prognum) at (\prognum*360/5+18 : 1.4cm and 1.2cm)
                {\includegraphics[width=.8cm]{icon-coding}};
        }
    }

    \opacityat{sm}{
        \node[image] (sm) {\includegraphics[width=.8cm]{icon-trainer}};
    }

    \node[
        draw=fg, thick, ellipse,
        inner sep=0cm,
        fit=(prog-1) (prog-2) (prog-3) (prog-4) (prog-5) (sm)
    ] (team) {};

    % Interface du PO
    \node[
        rectangle, fill=white,
        inner xsep=.6cm, inner ysep=.6cm
    ] at (sm-|team.east) {};

    \opacityat{po}{
        \node[image] (po) at (sm-|team.east)
            {\includegraphics[width=.8cm]{icon-user}};
    }

    % Parties prenantes
    \opacityat{pp}{
        \node[image, right=.2cm of po] (pp2)
            {\includegraphics[width=.8cm]{icon-woman}};

        \node[image, above=0cm of pp2] (pp1)
            {\includegraphics[width=.8cm]{icon-man}};

        \node[image, below=0cm of pp2] (pp3)
            {\includegraphics[width=.8cm]{icon-man}};
    }

    % Experts
    \opacityat{exp}{
        \node [image, left=.2cm of sm-|team.west] (exp2)
            {\includegraphics[width=.8cm]{icon-scientist}};

        \node [image, above=0cm of exp2] (exp1)
            {\includegraphics[width=.8cm]{icon-scientist}};

        \node [image, below=0cm of exp2] (exp3)
            {\includegraphics[width=.8cm]{icon-scientist}};
    }
\end{tikzpicture}

    \end{column}
\end{columns}
\end{frame}

\begin{frame}{Gestion du travail avec Scrum}
    \begin{columns}
        \begin{column}{.69\textwidth}
            \begin{itemize}
                \item \textbf{Backlog :} outil de gestion dynamique des tâches.
                \item Représente le produit sous forme de tâches priorisées.
                \item Unique, public et visible par tous.
                \item Évolue au cours du temps.
            \end{itemize}
        \end{column}

        \begin{column}{.31\textwidth}
            \includegraphics[width=\linewidth]{backlog}

            \hfill{\scriptsize \textbf{Source :} Logan Ingalls (CC BY 2.0)}\\[-.3em]
            \hfill{\scriptsize\url{https://flic.kr/p/91QewP}}
        \end{column}
    \end{columns}
\end{frame}

\chapter*{Conclusion et remerciements}
\addcontentsline{toc}{chapter}{Conclusion et remerciements}
\markboth{Conclusion}{Conclusion}
\label{chap:conclusion}

Ce projet nous a tout d'abord permis d'appliquer nos connaissances sur un problème concret et ainsi d' enrichir notre expérience. En effet, le fait de réaliser un projet de cette envergure nous a servi à nous tester sur nos compétences en programmation. Nous avions eu, d'ores et déjà au premier semestre, un défi de ce genre avec le jeu TenTen. Seulement, celui-ci, ne proposait pas la réflexion que nous a posé le Hashi puisqu'il s'agissait simplement d'un jeu et non d'un casse-tête. C'est pourquoi, avoir eu l'opportunité de travailler sur des projets comme celui-ci a pu améliorer notre expérience général, que ce soit en terme de programmation mais aussi de réflexion.\newline
Il nous a aussi entraîné à améliorer la gestion de notre temps, qui n'a pas toujours été optimale. Effectivement, afin de produire dans les temps un projet, il nous faut de l'organisation. Ce type "d'épreuve" permet ainsi de se prendre en main en nous faisant réaliser l'ampleur de ce qu'il y a faire. Cela nous aide à être encore plus autonome dans notre travail mais, pour autant, cela nous aide aussi à apprendre à travailler en groupe. Ces dernières notions sont d'autant plus importantes lorsqu'il s'agira de réaliser la demande de notre employeur ou de notre client quand nous trouverons un boulot après nos études. Ainsi, nous pouvons dire que ce type de projet est très important pour nous afin d'être au plus prêt lorsque cela arrivera.
\smallbreak
Malgré tout, notre programme n'est pas complet et ne permet pas de résoudre les problèmes les plus complexes, qui sont les cas où plus aucune de nos règles ne peuvent être appliquées. En effet, dans certaines configurations de départ, il se peut qu'aucune règle ne puisse être mis en oeuvre. Dans ce cas là, il faut alors commencer par essayer différentes combinaisons de ponts entre les îles de manière aléatoire et ainsi débloquer la situation, entraînant l'activation d'une des règles. Néanmoins, il est possible que la façon de commencer ne soit pas la bonne et mène à nulle part. Il faut alors repartir de zéro et, à nouveau, recommencer d'une autre manière jusqu'à trouver l'unique et bonne solution.\newline
Nous n'avons pas aussi eu le temps de s'attarder sur les règles concernant la connexité. Effectivement, il en existe et servent éventuellement à débloquer certaines situations ou du moins à les compléter. Dans notre version du Hashi, nous nous sommes penchés sur la connexité finale, c'est-à-dire, si nous avons bien qu'une seule composante connexe à la fin mais nous n'avons créé aucune règle similaire à celles existantes les concernant. Or il s'avère qu'elles peuvent permettre la création de ponts. Malgré cela, ces dernières ne sont pas indispensables, du moins, elles n'apparaissent que dans de rares cas.
\smallbreak
Pour finir, nous tenons particulièrement à remercier notre encadrant M. Philippe Janssen qui a su nous aider tout au long du projet grâce aux réunions hebdomadaires et notamment à l'aide de ses remarques sur notre travail mais aussi à l'aide de ses remarques sur notre organisation de travail. Cette aide nous a permis de nous débloquer à maintes reprises et nous a permis de savoir quoi faire lors de certaines situations.\newline


\end{document}
